\documentclass{article}
\usepackage{amsmath}
\usepackage{amssymb}
\usepackage{fullpage}
\usepackage{enumerate}
\usepackage{xcolor}
\pagestyle{empty}
\usepackage{graphicx}
\graphicspath{{C:\Users\steve\Pictures}}
\begin{document}

\noindent{{\Large \bf Fall 2025, 4100 Problem Set 3}

\large

\bigskip
\noindent{Please complete all problems. For any proof questions, please structure your proof similar to those from lecture.

\bigskip

\begin{enumerate}

\item Given that $\textbf{b}_1$ and $\textbf{b}_2$ are two nonzero, linear independent vectors in $\mathbb{R}^3$ and that $T\in\mathcal{L}(\mathbb{R}^4,\mathbb{R}^3)$ is such that 
$$T\textbf{e}_1=\textbf{b}_1,~~T\textbf{e}_2=-\textbf{b}_1,~~T\textbf{e}_3=2\textbf{b}_1,~~{\rm and}~~T\textbf{e}_4=\textbf{b}_2$$
find the general solution to $T{\bf x}={\bf b}_1$


\item Let $V$ be the vector space $\mathbb{P}_3(\mathbb{R})$ endowed with the inner product, $\langle x^i,x^j\rangle=\begin{cases}
1~{\rm if}~i=j\\
0~{\rm if}~i\neq j
\end{cases}$ and let $S$ be a subspace of $V$ given by, $S=\{f\in\mathbb{P}_3: f(1)=f'(1)=0\}$. Find bases for $S$ and $S^{\perp}$. 

\bigskip

{\large\bf Bonus}: Find an integral formula for the above inner product

\item Let $A\in\mathbb{R}^{m\times n}$ and $B\in\mathbb{R}^{m\times l}$. Show that ${\rm range}(A)$ is orthogonal to ${\rm range}(B)$ if and only if $A^TB$ is the zero matrix.


\item Compute the fundemental subspaces of the following matrix

$$A=\left[
\begin{array}{ccccc}
1 & 0 & 4 & -3 & 2\\
4 & -2  & 1 & -12 & 3\\
-3 & 2 & 3 & 9 & -1
\end{array}
\right]$$

\item Let 
$$A= \left[
\begin{array}{ccccccc}
1 & -1 & 0 & 0 & 0 & 0 & 0\\
-1 & 0 & 0 & 1 & 0 &  0 & 0\\
0 & 1 & 0 & -1 & 0 &  0 & 0\\
0 & -1 & 0 & 0 & 0 & 0  & 1 \\
0  & 0  & -1  & 0  & 1  & 0  & 0 \\
0  & 0  & 1  & 0  &  0 & -1  & 0\\
0 & 0 & 0 & 0 & -1 & 1 & 0
\end{array} \right] $$

\begin{enumerate}[a.]
\item Draw a directed graph having $A$ as its incidence matrix
\item Find bases for the four fundamental subspaces; ${\rm range}(A)$, ${\rm range}(A^T)$, ${\rm null}(A)$, and ${\rm null}(A^T)$
\end{enumerate}


\item Let $T$ be the linear map on $\mathbb{R}^2$ defined as follows
\begin{eqnarray*}
T\left[
\begin{array}{cc}
1 \\
0
\end{array}\right]&=\left[
\begin{array}{cc}
-2 \\
1
\end{array}\right]\\
T\left[
\begin{array}{cc}
2 \\
1
\end{array}\right]&=\left[
\begin{array}{cc}
0 \\
4
\end{array}\right]
\end{eqnarray*}
\begin{enumerate}[a.]
\item Find $[T]_{\{\textbf{e}_i\}}$ the matrix representation of $T$ with respect to the standard basis.
\item Let $\{\textbf{v}_1,\textbf{v}_2\}$ be another basis for $\mathbb{R}^2$ where $\textbf{v}_1=\left[
\begin{array}{ccc}
1 \\
-1 
\end{array} \right]$ and 
$\textbf{v}_2=
\left[
\begin{array}{ccc}
2 \\
-3 
\end{array} \right]$. Find $[T]_{\{\textbf{v}_i\}}$, the matrix representation of $T$ with respect to the basis $\{\textbf{v}_i\}$ and use it write 


$T(4{\bf v}_1+5{\bf v}_2)$ as a linear combination of the basis vectors, $\textbf{v}_i$.
\end{enumerate}


\item Let $T:\mathbb{P}_2\rightarrow\mathbb{R}_3$ be defined by $T(f)=f(-2){\bf e}_1+f(1){\bf e_2}+f(3){\bf e_3}$. Find the matrix representation of $T$ with respect to the standard bases of each space and use it to find the equation of the parabola passing through the points, $(-2,8),(1,3),(3,-5)$. 







\end{enumerate}
\end{document}